\lecture{4}{Tue 22 Sep 2026 09:31}{}

I missed 2 classes. Apparently we are now up to Van Kampen's theorem.

\begin{theorem}[Van Kampen]\label{thm:van-kampen}
	Let $X$ be the union of path-connected opens $A_\alpha $ each containing the basepoint $x\in X$, such that each $A_\alpha \cap A_\beta $ is path-connected. Then the map from the free product $\Phi : *_\alpha A_\alpha  \to X$ induced by the inclusion maps is surjective. If each triple intersection $A_\alpha \cap A_\beta \cap A_\gamma $ is also path-connected, Then the kernel is generated by $i_{\alpha \beta } (\omega )i_{\beta \alpha } (\omega )^{-1} $ for $\omega \in\pi _1(A_\alpha \cap A_\beta )$, and $i_{\alpha \beta } : \pi _1(A_\alpha \cap A_\beta ) \hookrightarrow \pi _1(A_\alpha)$ induced by the inclusion.
\end{theorem}


\begin{construction}
	A CW complex is similar to the realization of a simplicial set. Start with an $\mathbb{N} $-graded set $X^{\bullet} =\left\{ X^n \right\}_{n\geq 0} $. We think of an element $e^n\in X_n$ as an open $n$-ball $B^n$, and call it an \emph{$n$-cell}.

	Start by realizing $X^0$ as a discrete set of points. We write the realization of $X^0$ as $\left| X^0 \right| $. For every 1-cell $e^1_\alpha\in X^1$, we require a map $\varphi_\alpha^1 : \partial B^1 \to \left| X^0 \right| $. We then construct $\left| X^1 \right| $, the realization of all $X^n$ with $n\leq 1$, by gluing the 1-cells in along these maps:
	\[
		\left| X^1 \right| = \left(\left| X^0 \right| \amalg \coprod_{\alpha } B_\alpha ^1 \right)/ \left\langle \forall x\in \partial B^1_\alpha ,x \sim \varphi _\alpha ^1(x) \right\rangle 
	.\]
	We continue this process inductively:
	\[
		\left| X^n \right| =\left( \left| X^{n-1}  \right| \amalg \coprod_{\alpha } B^n_\alpha  \right) / \left\langle \forall x\in \partial B_\alpha ^n, x \sim \varphi_\alpha ^n(x) \right\rangle 
	.\]
	If $X^{\bullet} $ has only finitely many nonempty sets, then this is sufficient, as the realization $X = \left| X^{\bullet}  \right| $ is just $\left| X^n \right| $ for some $n$. However, if it is infinite, then we note that there are inclusions $\left| X^n \right| \subseteq \left| X^{n+1}  \right| $ for all $n$, and define $X$ as the colimit over these inclusions.

	We typically denote $\left| X^n \right| $ by $X^n$, and call it the $n$-skeleton. I only separated $X^n$ and $\left| X^n \right| $ for simplicity.
\end{construction}

Note that CW complexes are Hausdorff, locally contractible, and locally path connected.
